\section*{Capítulo \ref{ch:g2:measure} --- Dinheiro e Medidas}

\begin{solution}{exo:g2:measure:1}
Lápis: cm. Comprimento da sala de aula: m. Pena: g. Mochila escolar: kg. Água
garrafa: L.
\end{solution}

\begin{solution}{exo:g2:measure:2}
$2$ eu $= 200$ cm; \quad $300$ cm $= 3$ m; \quad $1$ kg $= 1000$ g.
\end{solution}

\begin{solution}{exo:g2:measure:3}
Por exemplo: $16 = 10 + 5 + 1$; \ $23 = 20 + 2 + 1$; \
$35 = 20 + 10 + 5$. (Outras maneiras são possíveis.)
\end{solution}

\begin{solution}{exo:g2:measure:4}
$10 - 6 = 4$; \quad $20 - 17 = 3$; \quad $50 - 32 = 18$.
\end{solution}

\begin{solution}{exo:g2:measure:5}
Mão longa $6$, curto entre $4$ e $5$: quatro e meia
($4{:}30$). Mão longa $3$, logo depois $8$: um quarto e meio
oito ($8{:}15$).
\end{solution}

\begin{solution}{exo:g2:measure:6}
Duas e meia: o ponteiro longo aponta para baixo, para o $6$, o ponteiro abreviado
fica entre o $2$ e o $3$.
\end{solution}

\begin{solution}{exo:g2:measure:7}
Depois de maio: junho. Antes de janeiro: dezembro. (O mês do aniversário varia.)
\end{solution}

\begin{solution}{exo:g2:measure:8}
Gasto: $3 + 4 = 7$. Mudar: $10 - 7 = 3$.
\end{solution}

\begin{solution}{exo:g2:measure:9}
$1$ kg $= 1000$ g, que é mais do que $800$ g: o açúcar é
mais pesado, por $1000 - 800 = 200$ g.
\end{solution}

\begin{solution}{exo:g2:measure:10}
Por exemplo $20 + 10 + 5 + 5 = 40$ e $10 + 10 + 10 + 10 = 40$
(duas vias com exatamente quatro moedas ou notas).
\end{solution}
